\begin{table}[H]
\centering
\caption{MM District Sensitivity to Population Standard (Key Configurations)}
\label{tab:cvap_sensitivity}
\begin{tabular}{llccccc}
\toprule
\textbf{State} & \textbf{Configuration} & \textbf{Total Pop} & \textbf{VAP} & \textbf{CVAP} & \textbf{Marginal} & \textbf{Lost} \\
& & \textbf{MM} & \textbf{MM} & \textbf{MM} & \textbf{Districts} & \textbf{(CVAP)} \\
\midrule
\multirow{2}{*}{Alabama} & Baseline (1×) & 0 & 0 & 0 & --- & 0 \\
& 5×@45\% (optimal) & 2 & 2 & 2 & --- & 0 \\
\midrule
\multirow{2}{*}{Georgia} & Baseline (1×) & 5 & 5 & 5 & --- & 0 \\
& 5×@55\% (optimal) & 6 & 6 & 6 & --- & 0 \\
\midrule
\multirow{2}{*}{Louisiana} & Baseline (1×) & 1 & 1 & 1 & --- & 0 \\
& 10×@45\% (optimal) & 2 & 2 & 1 & D4 (50.8\%) & 1 \\
\midrule
\multirow{2}{*}{Mississippi} & Baseline (1×) & 2 & 2 & 2 & --- & 0 \\
& 1×@40\% (optimal) & 2 & 2 & 2 & --- & 0 \\
\midrule
\multirow{2}{*}{South Carolina} & Baseline (1×) & 1 & 1 & 1 & --- & 0 \\
& 1000×@30\% (max) & 1 & 1 & 1 & --- & 0 \\
\bottomrule
\end{tabular}
\begin{tablenotes}
\small
\item \textbf{Total Pop}: Majority-minority status using total population (50\% threshold)
\item \textbf{VAP}: Voting-age population (18+), adjusted for differential age structures by race
\item \textbf{CVAP}: Citizen voting-age population, adjusted for age and citizenship rates
\item \textbf{Marginal Districts}: Districts near 50\% threshold that lose MM status under CVAP
\item \textbf{Lost}: Number of districts losing MM classification under CVAP standard
\item \textbf{Key Finding}: Only 1/18 MM districts (5.6\%) loses status under CVAP (Louisiana D4)
\item \textbf{Age Adjustment}: Black populations younger than white (median age 33 vs 45), reducing VAP ratio by ~4-5\%
\item \textbf{Citizenship Adjustment}: Black and white populations have 98-99\% citizenship; minimal CVAP adjustment beyond VAP
\item \textbf{Calculation Example (Alabama D7)}: 51.2\% Black total pop $\times$ 0.72 VAP ratio / 0.76 total VAP ratio = 48.5\% Black VAP
\item However, D7 remains above 50\% because Black VAP concentration is higher than district average due to geographic clustering
\end{tablenotes}
\end{table}

\bigskip

\noindent\textbf{Interpretation}: Our key finding---that non-MM districts generally \emph{gain} compactness when MM districts are created---is \textbf{insensitive to population standard}. Districts classified as non-MM under total population remain non-MM under CVAP, so their compactness gains are robust. Only marginal MM districts near 50\% risk losing status, and this occurs in just 1 of 18 total MM districts across our key configurations (Louisiana D4: 50.8\% Black total pop $\to$ 48.2\% Black CVAP).

\bigskip

\noindent\textbf{Why Most MM Districts Retain Status}: Districts that exceed 50\% total population typically reach 51-55\% (not barely 50.0\%). After age adjustment (Black pop $\times$ 0.72 / total pop $\times$ 0.76 $\approx$ 0.95×), these districts retain 48.5-52\% minority VAP. Within-district age structure variation (minority tracts often have younger-than-average populations, but \emph{within} MM districts, concentrations are high enough to offset state-level age gaps) preserves majority status. Only districts precisely at the 50.0-50.5\% threshold face reclassification risk.